\section{Paso 8: corridas de producción}
\label{sec:paso8-produccion}
% Origen: experimento PeatonalCorridasApareadas, corridas_peatonales.csv y cargar_corridas.py (2026-11-10).
% Días simulados y parámetros del experimento: anexos D y G.

La producción usa el experimento \texttt{PeatonalCorridasApareadas} con la configuración de la
tabla~\ref{tab:produccion}; los parámetros del modelo están en el anexo~\ref{anx:parametros}. Las cuatro
corridas de una réplica comparten la semilla, el día sorteado y todas las fuentes de azar.

\begin{table}[H]
\centering
\caption{Configuración de las corridas de producción}
\label{tab:produccion}
\small
\begin{tabularx}{\textwidth}{@{}>{\raggedright\arraybackslash}p{4.2cm} X@{}}
\toprule
Elemento & Valor \\
\midrule
Réplicas & Diseñadas: $n_0 = 10$ días hábiles. Completadas: 6, semillas \num{20260923} a \num{20260928} \\
Corridas & 24 de las 40 previstas: corrida $j$ con semilla $20260923 + \lfloor j/4 \rfloor$ y escenario $j \bmod 4$ (\EO{}, \EI{}, \ET{}, \EIT{}) \\
Día de cada réplica & Sorteado con la semilla entre los 75 días hábiles de marzo a junio de 2026 con servicio normal del Roca \\
Trenes & Llegadas reales de ese día, con carga de Poisson y el total del Roca igual al 80~\% de las validaciones del día \\
Validación & \EO{} y \EI{}: triangular (1,8; 2,4; 3,5)~s. \ET{} y \EIT{}: triangular (1,5; 1,8; 2,5)~s \\
Demora y descarga & 10~s de apertura de puertas; descenso por 21 puertas en 30~s \\
Franja & Ingresos de 07:00 a 09:30; la corrida sigue hasta que sale el último pasajero (tope técnico en \num{17900}~s) \\
Cierre por atasco & Pasadas las 09:30, si en 15~min no sale nadie y quedan a lo sumo $\max(20;\ 0{,}5~\%)$ pasajeros, se excluyen de las medidas y se informan (sección~\ref{sec:paso5-pilotos}) \\
Ruteo peatonal & Camino más corto, sin peso de densidad (elemento Pedestrian Settings) \\
Ejecución & Evaluaciones en paralelo, sin animación, AnyLogic~8.9.10 PLE \\
\bottomrule
\end{tabularx}
\end{table}

\textbf{Controles.} El cargador \texttt{cargar\_corridas.py} leyó 24 filas \texttt{CSV\_PEATONAL} (6 semillas con
\EO{} y 6 pares completos por alternativa), ninguna \texttt{INCOMPLETO}. En cada corrida se cumplió la
conservación, con entre 19 y 66 atascados (0,1 a 0,4~\% de los pasajeros, por debajo del tope de descarte), y en
cada réplica los cuatro escenarios tuvieron el mismo día y los mismos generados. Ninguna fila se transcribió a
mano. Los intervalos salen de \texttt{analisis\_corridas.py}, porque la planilla
\texttt{04-resultados-corridas.xlsx} los habilita recién con 10 pares; \texttt{resultados\_informe.py} genera
las tablas y figuras con las mismas funciones. Los seis días sorteados (del 06/03 al 05/06/2026, con 43 a 50
trenes y entre \num{16760} y \num{20530} pasajeros generados) se listan en el anexo~\ref{anx:dias}. Cada corrida
mueve unos \num{18000} peatones y una tanda de 10 en paralelo tardó del orden de 100 minutos: por eso se
completaron 6 de los 10 días del diseño antes de la entrega.
